\section{Introduction}

Belief elicitation in experiments often asks subjects to report probabilities directly.
Frequency reports ask instead for counts: out of \(n\) realizations, how many times will each outcome occur?
This format is natural when the object of interest is a distribution of choices, signals, or outcomes in a finite group.
It may also be easier for subjects than reporting abstract probabilities under a standard scoring rule.
Many probability scoring rules, including the quadratic scoring rule, are well understood theoretically \cite{Savage1971,GneitingRaftery2007,Selten1998}, but they require subjects to understand how probability reports map into payments.
That mapping can be cognitively demanding, especially when subjects reason more naturally in frequencies than in abstract probabilities \cite{Hogarth1975,SchlagEtAl2013}.
Direct monetary scoring rules also raise the familiar risk-aversion problem \cite{ArmantierTreich2012}.

The practical problem is that the researcher observes a count report \(r\), not the subject's latent belief vector \(p\).
The report is an incentivized message generated by a scoring rule.
Thus the inferential question is inverse: given the rule and the observed report, which beliefs are consistent with optimal reporting?
The paper follows the chain
\[
\text{frequency-report scoring rule}
\rightarrow
R_S(p)
\rightarrow
P_S(r)
\rightarrow
\text{belief and functional bounds}.
\]
Here \(R_S(p)\) is the set of reports that are optimal for a subject with beliefs \(p\), and
\[
P_S(r)=\{p\in\Delta^k:r\in R_S(p)\}
\]
is the inverse belief region induced by the observed report.
From this set, the researcher can derive bounds on individual probabilities and on functionals such as expected values.

This paper uses the term informational efficiency in an operational sense.
A rule is more informative for a given design objective if the observed report leads to narrower coordinate bounds or narrower mean and functional bounds.
These inferential gains must be weighed against implementation costs: exact-match payment probabilities, cognitive simplicity, computational burden, and robustness to risk aversion.

The paper's main analytical contribution is for quadratic-distance frequency scoring.
Under this rule, a subject reports the feasible count vector closest to the expected count vector \(np\).
This gives a linear belief region and closed-form bounds for each \(p_i\).
Mean bounds and other linear functional bounds are then simple linear programs over the same region.

The paper compares this rule to several other count-loss mechanisms.
Schlag--Tremewan frequency guessing is the known discrete-metric exact-match mechanism: it elicits multinomial modes, gives closed-form bounds, and is robust to risk aversion under a fixed-prize implementation \cite{SchlagTremewanSimple}.
Manhattan-distance scoring is included as a median-oriented comparison with an explicit threshold representation.
Hamming and Chebyshev rules are discussed as secondary computable rules because their inverse regions are exactly specified by finite expected-loss inequalities, but they are not the focus of the headline design exercise.

The contribution is therefore not a new exact-match mechanism and not a general theory of scoring rules.
It is a finite-sample method for interpreting frequency reports as belief and mean bounds, together with a design comparison of simple mechanisms that differ in transparency, robustness, and inferential precision.
